% Stronger section fragment.  Reuses notation and tree-copy construction in
% exceptional_columns_theorems.tex.  This supersedes its factorial FPT bound.
\section{Compressing all exceptional schedules into a Boolean determinant filter}
\label{sec:special-row-sieve}

The schedule method exposes a deterministic exact algorithm, but its
factorial enumeration is unnecessary for randomized discovery.  Every
exceptional column can change only one of its incident rows.  Keeping track
of those rows captures all nonlinear activation, while ordinary rows are
handled by graph reachability and the fixed source-component constraints.
This gives a single-exponential parameter bound.

Let $P$ and $Z$ be the ordinary implication graph and ordinary seed set,
and let
\[
             U=\bigcup_{c\in E}\operatorname{supp}(c),
             \qquad t=|U|\leq3k.
\]
Use $\operatorname{Cl}(X)$ for the closure of $X\cup Z$ under ordinary
reachability and all exceptional rules.  For a triple $\{u,v,w\}$ these
rules are $\{u,v\}\to w$, $\{u,w\}\to v$, and $\{v,w\}\to u$.
Equivalently, $\operatorname{Cl}(X)$ consists of the rows removed by
peeling after initially deleting $X$.  All rules concern removed rows.

Define a monotone Boolean function on subsets of the special rows by
\begin{equation}\label{eq:special-truth}
             f(A)=\mathbf 1\{U\subseteq\operatorname{Cl}(A)\},
                    \qquad A\subseteq U.
\end{equation}
The seeds $Z$ are always included in this definition.  They will not be
encoded as colors that a tree is required to hit.  In particular $f$ may
already be one on the empty set.

Let $S_1,\ldots,S_q$ now denote the source SCCs of the ordinary graph $P$
that are disjoint from $Z\cup U$.  These groups are fixed for the entire
algorithm, independently of the tree and the Boolean input $A$.
For a spanning tree $T$, define its special influence set by
\begin{equation}\label{eq:special-influence}
        H(T)=\operatorname{Reach}_P(T)\cap U
             =\{u\in U:T\cap\operatorname{Anc}(u)\ne\varnothing\}.
\end{equation}
Only influence from $T$ is included here.  Influence from the free seeds
is already accounted for in $f$.

\begin{lemma}[Exact special-row criterion]\label{lem:special-criterion}
A spanning tree $T$ leaves a completely peelable matrix if and only if
\begin{equation}\label{eq:special-criterion}
       T\cap S_a\ne\varnothing\quad(1\leq a\leq q),
                    \qquad f(H(T))=1.
\end{equation}
\end{lemma}
\begin{proof}
First we prove that the special rows activated from $T$ are exactly those
activated from $H(T)$.  Before any exceptional rule fires, their special
sets are respectively
\[
     \operatorname{Reach}_P(Z\cup T)\cap U
       \quad\text{and}\quad
     \operatorname{Reach}_P(Z\cup H(T))\cap U.
\]
These sets agree.  Every member of $H(T)$ is reachable from $T$; conversely
every special row reachable from $T$ is already a member of $H(T)$.
If the special sets agree at some stage, precisely the same exceptional
rules are enabled, because every premise and head of every such rule
lies in $U$.  Their new heads also lie in $U$.  Ordinary closure from a
common set of new special heads has the same effect on special rows.
Induction over the finite monotone closure process proves
\begin{equation}\label{eq:special-closure-equality}
          \operatorname{Cl}(T)\cap U
             =\operatorname{Cl}(H(T))\cap U.
\end{equation}
There is no need to retain the initially activated ordinary rows outside
$U$: their only future effect on special rows has already been included
in $H(T)$.

If all rows peel, \eqref{eq:special-closure-equality} gives $f(H(T))=1$.
A source SCC disjoint from $Z\cup U$ cannot receive an ordinary implication
from another SCC and cannot contain an exceptional head or a seed.
It therefore must contain a deleted tree row.  This proves necessity.

Conversely suppose \eqref{eq:special-criterion}.  All special rows become
active by \eqref{eq:special-closure-equality}.  Every source SCC of $P$
then contains an active row: a seed in $Z$, a special row in $U$, or a tree
row by the hitting condition.  Every SCC of a finite condensation DAG is
reachable from a source SCC.  Ordinary closure therefore activates all
rows.  This is complete peeling.
\end{proof}

\subsection{Computing the activation table without expanded words}

The table of $f$ can be computed in $2^t\operatorname{poly}(N)$ time.
Precompute ordinary reachability between special rows and the free mask
$\operatorname{Reach}_P(Z)\cap U$.  For each initial mask $A$, repeatedly
add the ordinary successors of active special rows and every exceptional
head whose prerequisites are active.  At most $t$ strict mask enlargements
can occur.  The resulting mask is all of $U$ exactly when $f(A)=1$.
This is a finite monotone closure computation on bit masks, not an
exponential expansion in normal-surface weight.

For a more economical parameter, let $d$ be the number of ordinary SCCs
that meet $U$.  Two rows of one SCC have identical ordinary ancestor sets,
and one becomes active ordinarily whenever the other does.  Hence one
color per such SCC suffices.  Choose one representative of each SCC,
interpret an initial color mask as activating its representatives, and
use the same closure test on all special rows.  Thus all statements below
hold with $t$ replaced by
\begin{equation}\label{eq:activation-dimension}
             d\leq t\leq3k.
\end{equation}
For clarity the algebraic proof is written with the individual-row set
$U$; replacing it by SCC colors changes none of its identities.

\subsection{An arbitrary Boolean filter on common bases}

Use the copied graphic-incidence/Vandermonde construction from
Lemma~\ref{lem:tree-cover-sieve}, with the fixed groups $S_1,\ldots,S_q$
from Lemma~\ref{lem:special-criterion}.  If $q>r$, no tree can hit all
these disjoint groups.  Otherwise common bases of the two represented
matroids correspond exactly to spanning trees that hit all these groups.
There are at most $M\leq2m$ copies, each with its own variable $z_e$.

Assign to a copy of a row $e$ the color mask
\[
             C(e)=\{u\in U:e\in\operatorname{Anc}(u)\}.
\]
The union of the masks in a copied spanning-tree basis is exactly $H(T)$.
For $A\subseteq U$, let
\[
     P_A(z)=\det\bigl(A_{\rm gr} D_A(z)B^{\mathsf T}\bigr),
\]
where $D_A(z)$ retains $z_e$ if $C(e)\cap A=\varnothing$, and sets that
entry to zero otherwise.  Thus a common basis $F$ occurs in $P_A$ exactly
when its hit mask $H(F)$ avoids $A$.

Define
\begin{equation}\label{eq:boolean-mobius}
     g(K)=f(U\setminus K),\qquad
     c_A=\sum_{B\subseteq A}(-1)^{|A\setminus B|}g(B),
\end{equation}
and form the polynomial
\begin{equation}\label{eq:boolean-determinant-filter}
                        Q(z)=\sum_{A\subseteq U}c_A P_A(z).
\end{equation}

\begin{theorem}[Boolean common-basis sieve]\label{thm:boolean-basis-sieve}
The polynomial $Q$ is multilinear and homogeneous of degree $r$.  Its
monomials with nonzero coefficient encode exactly the copied spanning-tree
bases whose underlying trees satisfy \eqref{eq:special-criterion}.
In particular
\[
              Q\not\equiv0
                \quad\Longleftrightarrow\quad
              \text{forest-compatible raw completion exists}.
\]
This holds over any field for which the two matrix representations are
valid.  The Boolean function need not be monotone for the algebraic
identity; monotonicity enters only its interpretation as activation.
\end{theorem}
\begin{proof}
Cauchy--Binet associates a common basis $F$ with the distinct monomial
$\prod_{e\in F}z_e$ and its nonzero coefficient
$\det((A_{\rm gr})_F)\det(B_F)$.  It occurs in $P_A$ precisely when
$A\subseteq U\setminus H(F)$.  Its coefficient in $Q$ is therefore its
original coefficient multiplied by
\[
          \sum_{A\subseteq U\setminus H(F)}c_A
             =g(U\setminus H(F))=f(H(F)).
\]
The middle identity is subset-lattice M\"obius inversion.  To verify it
directly, substitute \eqref{eq:boolean-mobius}, interchange the two finite
sums, and use
$\sum_{A:\,B\subseteq A\subseteq K}(-1)^{|A\setminus B|}=0$
unless $B=K$, in which case it is one.  Thus the multiplier is exactly
zero or one.  The identity holds over the integers and after reduction
in any characteristic.  Large or negative intermediate $c_A$ do not
introduce an additional characteristic restriction.

Distinct bases use distinct copy variables, so different encodings do
not share a monomial.  Every surviving basis coefficient remains nonzero.
The exact special-row criterion now proves the final equivalence.
Multilinearity and homogeneity follow directly from the degree-$r$
Cauchy--Binet monomials.
\end{proof}

\subsection{Single-exponential discovery with checked witnesses}

\begin{theorem}[Activation-dimension algorithm]\label{thm:activation-fpt}
Let $d$ be the number of ordinary SCCs incident with exceptional columns.
For every $0<\varepsilon<1$, there is a randomized algorithm with time
\begin{equation}\label{eq:activation-time}
                 2^d\operatorname{poly}(N,\log(1/\varepsilon))
                      \ \leq\ 
                 8^k\operatorname{poly}(N,\log(1/\varepsilon))
\end{equation}
that finds a spanning tree and a verified unit-column peel order with
probability at least $1-\varepsilon$ whenever a completion exists.
It never reports an incorrect positive certificate.  It uses
$O(2^d\log N+\operatorname{poly}(N))$ bits of space.  There is also a
polynomial-space variant with time
\[
                   3^d\operatorname{poly}(N,\log(1/\varepsilon)).
\]
\end{theorem}
\begin{proof}
Compute ordinary SCCs, ancestor masks, and the truth table on the $d$
SCC colors.  Choose a prime field with characteristic
$P>\max(M,2r)$, of polynomial magnitude in $N$.  The represented graphic
matroid and the filler Vandermonde matroid are valid over this field.
Compute $g$ and its subset M\"obius transform $c$ by the standard in-place
recurrence
\[
     c[A]\leftarrow c[A]-c[A\setminus\{i\}]
       \quad\text{for each color }i\text{ and each }A\ni i.
\]
There are $O(d2^d)$ field operations.  All coefficients may be computed
modulo $P$; their integer identities have already been proved.

An evaluation of $Q$ uses $2^d$ determinants of order $r$ and polynomial
additional work.  If completion exists, $Q$ is a nonzero degree-$r$
polynomial, so a uniform random assignment misses with probability at
most $r/P<1/2$.  With
$b=\max\{1,\lceil\log_2(1/\varepsilon)\rceil\}$ independent trials,
the probability of missing every time is at most $\varepsilon$.
A nonzero evaluation guarantees existence without qualification.

Use the deterministic repeated-deletion extraction from
Lemma~\ref{lem:tree-cover-sieve}, now with $Q$ in place of $\Phi$.
All of its hypotheses hold: $Q$ is homogeneous and multilinear of degree
$r$, and $P>M$.  At most $O(M^2)$ further evaluations produce exactly
$r$ copies forming a common basis with $f(H)=1$.  Decode its tree and
run the independent unit-column peeling checker on the original matrix.
Lemma~\ref{lem:special-criterion} proves that every row is removed.
No random choices are used during extraction, so the error bound remains
that of the initial identity tests.

Storing $f$, $g$, and $c$ costs $O(2^d\log N)$ bits when values are stored
modulo $P$, in addition to the polynomial graph and matrix storage.
For polynomial space, do not store these tables.  Compute each $c_A$
directly from its subsets, evaluating $f$ by closure as needed.  Summed
over all $A$, the number of truth-value computations is
\[
                          \sum_{A\subseteq[d]}2^{|A|}=3^d.
\]
Every determinant, coefficient, and restriction can then be processed
sequentially.  The repeated extraction calls multiply the time only by
a polynomial factor.
\end{proof}

The randomized algorithm has one-sided misses.  A production interface
can report \texttt{SUCCESS} only after exact replay, and return
\texttt{INCONCLUSIVE} when the tested assignments are all zero.  An exact
negative decision may use the deterministic schedule solver, an existing
stopping-set certificate, or exhaustive search.  Probabilistic failure
to find a witness is not a topological obstruction.

\begin{corollary}[Larger exceptional supports]
The activation-dimension theorem remains valid if exceptional active
columns have arbitrary finite support and arbitrary positive counts,
provided that $U$ contains every row incident with them.  For a unit entry
on row $h$, its rule requires every other support row to have been removed.
Columns with no unit entry remain inert.  The running time is
$2^d\operatorname{poly}(N,\log(1/\varepsilon))$, where $d$ counts the
ordinary SCCs meeting this larger exceptional-row set.
\end{corollary}
\begin{proof}
Every exceptional rule still has all its premises and head in $U$.
The exact special-row criterion and Boolean filter use no bound on rule
arity.  A finite closure test is polynomial in the explicitly encoded
rule incidence.  The triple-specific inequality $d\leq3k$ is the only
part that no longer applies.
\end{proof}

\paragraph{Concrete consequence and its boundary.}
If current states have encoded size $N=n^{O(1)}$ and activation dimension
$d=O(\log n)$, this local tree-selection problem has randomized
polynomial-time certificate discovery.  For $d=O((\log n)^a)$ it has
quasi-polynomial local time.  No such dimension bound is established for
all states of a general unknot hierarchy.  The theorem also leaves intact
fixed-interface stopping obstructions: it decides the same restricted
geometric move class more efficiently, and does not make that class
complete for handlebody recognition.
